\section{Appendix}

\subsection{Local setup and test execution}
To install, configure, and execute Rùa Toán Basic locally from the repository:
\begin{lstlisting}[language=bash,numbers=none]
# 1. Clone repository and set up Python virtual environment
git clone https://github.com/thunopro/ruatoan_basic.git
cd ruatoan_basic
python3 -m venv .venv
source .venv/bin/activate
pip install -r requirements.txt

# 2. Configure environment variables (.env.local)
cp .env.example .env.local
# Edit .env.local: set DATABASE_URL (PostgreSQL) and SESSION_SECRET

# 3. Execute automated end-to-end smoke test suite (all 8 flows, 26 checks)
python scripts/smoke_test.py

# 4. Launch local development server
uvicorn main:app --reload --port 8000
# Open http://127.0.0.1:8000 in browser
\end{lstlisting}
On startup, the database module connects using \texttt{DATABASE\_URL}, sets up schema
\texttt{ruatoan\_basic}, and creates tables \texttt{users}, \texttt{photos}, and \texttt{password\_resets}.
Environment secrets are strictly excluded from source control.

\subsection{Demo accounts}
Table~\ref{tab:demo-accounts} provides pre-configured test accounts for the instructor and evaluators to test
authenticated features on the live deployment (\url{https://ruatoanbasic.vercel.app}).

\begin{table}[ht]
\centering\small
\begin{tabularx}{0.92\textwidth}{@{}l l l >{\raggedright\arraybackslash}X@{}}
\toprule
\textbf{Role} & \textbf{Email} & \textbf{Password} & \textbf{Testing capabilities}\\
\midrule
Student Evaluator & \texttt{demo@ruatoan.vn} & \texttt{StudentPass123!} & Problem bookmarking, homework photo upload/delete, credential change\\
Clean Test Account & \texttt{test@ruatoan.vn} & \texttt{TestUser123!} & Fresh account for registration flow and single-use password reset testing\\
\bottomrule
\end{tabularx}
\caption{Pre-configured demo accounts for instructor and grader evaluation.}
\label{tab:demo-accounts}
\end{table}

\subsection{Route inventory}
Table~\ref{tab:routes} provides a compact route map for the pages and API endpoints implemented in Rùa Toán Basic.

\begin{table}[ht]
\centering\small
\begin{tabularx}{\textwidth}{@{}l >{\raggedright\arraybackslash}X@{}}
\toprule
\textbf{Route} & \textbf{Purpose and responsibilities}\\
\midrule
\texttt{/} & Home dashboard: live catalogue metrics, mission, and marked problem bookmarks\\
\texttt{/pages/practice.html} & Problem search, grade/topic filters, pagination, and MathJax formula rendering\\
\texttt{/pages/exams.html} & Curated collection list and detailed collection view with ordered statements\\
\texttt{/pages/lesson.html} & Theory lessons with client-side keyword search and accordion LaTeX bodies\\
\texttt{/pages/library.html} & Categorized reference books and documents with search and PDF links\\
\texttt{/pages/login.html}, \texttt{/pages/register.html} & User authentication forms with credential validation\\
\texttt{/pages/submissions.html} & Student homework photo upload, thumbnail gallery, and deletion\\
\texttt{/css/*}, \texttt{/js/*}, \texttt{/images/*} & Static assets served directly by Vercel global edge CDN\\
\texttt{/api/stats} & Live catalogue totals and multi-topic counts\\
\texttt{/api/daily} & Deterministic daily practice problem\\
\texttt{/api/collections}, \texttt{/api/collections/\{id\}} & Collection list (summaries) and detail with pagination\\
\texttt{/api/problems} & Problem search with folded Vietnamese diacritics and pagination\\
\texttt{/api/lessons}, \texttt{/api/books}, \texttt{/api/docs} & Reference catalogue collections\\
\texttt{/api/register}, \texttt{/api/login}, \texttt{/api/logout}, \texttt{/api/me} & Account session management with signed cookies\\
\texttt{/api/change-password} & Credential update with current password validation\\
\texttt{/api/request-reset}, \texttt{/api/reset-password} & Database-persisted single-use password reset\\
\texttt{/api/photos}, \texttt{/api/photos/\{id\}} & Photo upload (magic bytes), list, binary stream, and delete\\
\bottomrule
\end{tabularx}
\caption{Application page and API route map.}
\label{tab:routes}
\end{table}

\subsection{Content inventory snapshot}
The audited curriculum data files contain the following verified records:
\begin{itemize}
  \item 11 source JSON files: 3,381 collections and 33,929 nested problem records.
  \item 7 theory JSON files: 90 lessons across five mathematical branches.
  \item \texttt{books.json}: 216 books, magazines, and resource records.
  \item \texttt{docs.json}: 89 document records.
\end{itemize}

\subsection{Data pipeline diagram}
Figure~\ref{fig:content-pipeline} illustrates how offline curriculum preparation feeds runtime consumption.

\begin{figure}[ht]
\centering
\begin{tikzpicture}[
  stage/.style={draw=PrussianBlue,fill=PaleBlue,rounded corners=2pt,align=center,text width=2.7cm,minimum height=1.2cm},
  output/.style={draw=FDABlue,fill=Solitude,rounded corners=2pt,align=center,text width=2.7cm,minimum height=1.2cm},
  arrow/.style={->,thick,color=PrussianBlue}
]
\node[stage] (source) at (0,0) {External sites and\\exam/magazine/\\book sources};
\node[stage] (collect) at (3.7,0) {Crawlers and\\intermediate cache\\(PDFs / metadata)};
\node[stage] (prepare) at (7.4,0) {Extract, clean,\\normalise, tag,\\retain attribution};
\node[output] (json) at (11.1,0) {Versioned JSON\\under \texttt{data/}};
\node[stage] (runtime) at (11.1,-2.1) {Content parser\\cache and in-memory\\lookup maps};
\node[output] (site) at (7.4,-2.1) {Plain HTML/CSS/JS\\browser views};
\draw[arrow] (source) -- (collect);
\draw[arrow] (collect) -- (prepare);
\draw[arrow] (prepare) -- (json);
\draw[arrow] (json) -- (runtime);
\draw[arrow] (runtime) -- (site);
\end{tikzpicture}
\caption{Content preparation and runtime consumption. The pipeline is an offline tool; the basic platform reuses its structured JSON output directly.}
\label{fig:content-pipeline}
\end{figure}

\subsection{Glossary}
\begin{description}
  \item[Collection] An ordered source-aware unit such as an exam paper or topic exercise set.
  \item[Problem] One individually addressable statement nested in a collection.
  \item[Topic tag] A classification label; one problem may have multiple tags.
  \item[Session cookie] An HMAC-signed \texttt{HttpOnly} cookie holding the user id and expiry.
  \item[Magic bytes] The unique initial bytes of a file format (e.g. \texttt{\textbackslash x89PNG}) used to reliably identify binary file types.
\end{description}

\subsection{AI assistance and academic disclosure}
In accordance with course guidelines and academic integrity standards, the group discloses that
AI assistance (Anthropic Claude and Google Gemini) was utilized during the development and report
preparation process:
\begin{itemize}
  \item \textbf{Ideation and architectural consultation:} Assisting in identifying the pedagogical limitations
        of the initial Next.js platform and formulating the foundational plain HTML/CSS/JS architecture.
  \item \textbf{Code debugging and hardening:} Assisting the team in writing defensive patterns, including the
        magic-byte validation function, diacritic normalization helper, and \texttt{psycopg} connection pool configuration.
  \item \textbf{Report drafting and formatting:} Assisting in formatting LaTeX document structures, generating
        TikZ diagrams, and structuring tables to adhere to the Faculty of Data Science and AI template.
\end{itemize}
All code, database schemas, and documentation were reviewed, comprehended, tested, and verified directly by the team.
